\section{Introduction}
\label{sec:intro}

% TODO: full draft. Suggested argument (mirrors the related-works arc):
%
% P1 — The problem. ~57M people live with dementia; informal caregivers
%   shoulder daily care with little training and elevated depression/
%   anxiety risk. They already seek answers online and from LLMs.
%   Cite: \citep{forumanalysis2024,balancing2025}.
%
% P2 — The gap. NLP for dementia focuses on detection, not caregiver
%   support \citep{nlpdementia2025}; the care community sees LLM promise
%   \citep{llmdementia2024}, but ungrounded LLMs omit critical details
%   60-80% of the time \citep{homecare2025} and caregivers distrust
%   uncited answers \citep{carey2025}.
%
% P3 — The opportunity. Decades of expert geriatric teaching exist as
%   long-form lecture video (e.g., dementia-care educator lectures) —
%   a vetted, mentorship-style knowledge source no prior system uses.
%
% P4 — Our approach. Extend MentorQA \citep{bhalerao2026mentorqa} to
%   healthcare: chunk lecture transcripts, generate mentorship-focused
%   QA with a multi-agent pipeline, ground answers in source segments,
%   evaluate with LLM judges + human pilot.
%
% P5 — Contributions (bulleted): (1) first mentorship-oriented QA
%   resource grounded in expert dementia-care lectures; (2) controlled
%   comparison of QA-generation architectures in a clinically dense
%   domain; (3) evaluation framework with LLM-judge/human agreement
%   analysis; (4) released dataset + code.

\todo{Write introduction following the outline in the comments above.}
